% #########################################################################
% #               METODOLOGÍAS Y TÉCNICAS DE INVESTIGACIÓN                #
% #########################################################################

% Definición de nombres
\newcommand{\estudiante}{García Justel, Alan}
\newcommand{\titulo}{MÁSTER EN INGENIERÍA COMPUTACIONAL Y SISTEMAS INTELIGENTES}
\newcommand{\asignatura}{METODOLOGÍAS Y TÉCNICAS DE INVESTIGACIÓN}
\newcommand{\portada}{common/no_signal.png}
\newcommand{\colorportada}{title_red}
\newcommand{\curso}{2024-2025}



% Notebook
\begin{document}

\newgeometry{bottom=2cm}

\begin{titlepage}
    % Logo de la universidad
    \begin{textblock*}{\textwidth}(2cm,0.4cm)
        \begin{center}
            \begin{minipage}{0.45\textwidth}
                \centering
                \includegraphics[width=\textwidth]{common/Logo_EHU.jpg}
            \end{minipage}\hfill
            \begin{minipage}{0.45\textwidth}
                \centering
                \includegraphics[width=0.8\textwidth]{common/Logo_KISA.png} 
            \end{minipage}
        \end{center}
    \end{textblock*}
    
    % Franja de color
    \begin{tikzpicture}[remember picture, overlay]
        \fill[\colorportada] (current page.north west) ++ (0,-3.01cm) rectangle (\paperwidth,-3cm);
    \end{tikzpicture}
    
    \begin{textblock*}{\paperwidth}(\dimexpr\parindent+\oddsidemargin+3em\relax,3.5cm)
        \begin{minipage}{\dimexpr\linewidth-7.5cm\relax}
            \color{white}
            \noindent\rule{\linewidth}{0cm}
            \textsf{ {\large \titulo}}
            \newline
            \newline \newline
            \textsf{\textbf{ {\Huge APUNTES DE ASIGNATURA }}}
        \end{minipage}
    \end{textblock*}
    
    % Nombre asignatura
    \vspace*{3.5cm}
    \begin{minipage}{\linewidth}
        \setlength{\baselineskip}{1.7\baselineskip}
        \centering
        \textsf{ \textbf{ {\LARGE \asignatura }}}
    \end{minipage}

    % Foto de portada
    \vspace*{0.5cm}
    \begin{figure}[H]
        \centering
        \includegraphics[width=10cm, height=8cm]{\portada}
    \end{figure}

    
    
    % Estudiante
    \vspace{0.2cm}
    \noindent {\footnotesize \textbf{Estudiante:} \estudiante}
    \newline
    \noindent\makebox[\linewidth]{\rule{\textwidth}{0.4pt}} % Línea horizontal

    % Curso y Fecha
    \vspace{0.1cm}    
    \noindent {\footnotesize \textbf{Curso: } \curso \hfill \textbf{Fecha:} \today }
\end{titlepage}

\restoregeometry
\setcounter{figure}{0} % Incluimos el título
\newpage

% Índices
% \tableofcontents\thispagestyle{empty} %\newpage
% \listoffigures\thispagestyle{empty}   %\newpage
% \listoftables\thispagestyle{empty}    %\newpage

\section{Introducción al mundo de la investigación}
Existen varios tipos de documentos científicos. Por un lado tenemos las tesis de doctorado o de máster en las que se sintetiza un trabajo de investigación que aporta nuevos resultados sobre el presente estado del arte del momento. Por otro lado, existen los artículos científicos que se presentan en revistas y congresos. Estos últimos aportan mucho <<Networking>> de tal forma que te permite saber las líneas de investigación existentes sobre un tema actual.

Finalmente, para poder investigar es necesario contar con una <<Solicitud de Proyecto de Investigación>>. Es un documento científico importante en el que además hay que hacer un análisis exahustivo del estado del arte y es necesario para conseguir fondos para el proyecto. 

\subsection{Empezar a escribir una tésis}
El primer paso es la elección y justificación del tema o, si ya se tiene elegido, búsquedas de documentos relacionados con el tema a tratar. Para ello, existen varios portales como \href{https://scholar.google.es/}{Google Scholar}, \href{https://www.educacion.gob.es/teseo/irGestionarConsulta.do}{TESEO}, \href{https://www.proquest.com/}{ProQuest}... Además, hay que justificar el problema y mostrar la utilidad del trabajo a realizar (parte del contexto) y un análisis del estado del arte: antecedentes, métodos actuales, limitaciones y condicionantes etc. Otro aspecto muy importante es el planteamiento de los objetivos de la tésis.

Antes de comenzar con la escritura, es muy recomendable desglosar el tema y dividirlo en capítulos. En resúmen, es recomendable realizar una exploración sobre las líneas actuales de investigación afines al tema a tratar y crear el índice del documento antes de comenzar con la escritura.

\ex{Estructura de una tésis}{
    \begin{enumerate}
        \item Introducción
        \item Contexto / Comparativa con el estado del arte
    \end{enumerate}
}


\nt{Papers interesantes}{
State of the art of semantic segmentation for autonomous driving:
    \begin{enumerate}
        \item \url{https://ieeexplore.ieee.org/abstract/document/8206396}
        \item \url{https://link.springer.com/chapter/10.1007/978-3-031-19769-7_23}
        \item \url{https://link.springer.com/article/10.1007/s00138-023-01400-7?fromPaywallRec=false}
        \item \url{https://link.springer.com/chapter/10.1007/978-981-99-8435-0_29?fromPaywallRec=false}
    \end{enumerate}
}

\subsection{Revisión del estado del arte}
La revisión del estado del arte es un proceso metódico para identificar, evaluar e interpretar el trabajo hecho por otros investigadores sobre un tema de interés. No es muy común el referenciar tésis doctorales o de máster durante esta revisión y, aunque no habría problema en referenciar un par de ellas, sí que es muy importante referenciar artículos de primer nivel para dar peso al análisis del estado del arte.

En definitiva, la sección de la revisión del estado del arte es el resultado de un proceso de análisis, síntesis y comprensión que se refleja a modo de una pequeña <<historia>>.

\qs{Búsqueda de grupos de investigación}{
Se tienen que buscar grupos de investigación de departamentos de informática actuales en la universidad de Granada (UGr) y recoger cuáles son sus líneas de investigación principales.
}

\sol{
    Los siguientes grupos de investigación se han obtenido del siguiente \href{https://lsi.ugr.es/investigacion/grupos#title0}{link} a fecha de $17/09/2024$:
    {\tiny
    \begin{enumerate}[itemsep=0mm]
        \item Grupo de Especificación, Desarrollo y Evolución de Software (GEDES)\begin{itemize}[nosep]
            \item Software para discapacidad
            \item Enseñanza Virtual
            \item Informática y accesibilidad
            \item Modelos epistemológicos
        \end{itemize}
        \item Informática Gráfica y Realidad Virtual
        \begin{itemize}[nosep]
            \item Realidad Virtual.
            \item Digitalización 3D de patrimonio histórico
            \item Representación y visualización de modelos 3D para aplicaciones médicas
            \item Animación por ordenador y visualización expresiva.
            \item Ingeniería de aplicaciones gráficas.
            \item Visualización de grandes modelos.
            \item Maquetación Virtual.
        \end{itemize}
        \item Modelling and Development of Advanced Software Systems (MYDASS)
        \begin{itemize}[nosep]
            \item Human Computer Interaction
            \item Assistive technologies
            \item Frameworks and platforms for mobile systems
            \item Sensor Networks and domotics for Health Systems
        \end{itemize}
        \item Sistemas Concurrentes
        \begin{itemize}[nosep]
            \item Verificación de software asistida por computador
            \item Métodos formales en sistemas concurrentes
            \item Computación de tiempo real y empotrada para sistemas ubicuos y de inteligencia ambiental
            \item Middlewares y marcos de desarrollo para sistemas empotrados de tiempo real distribuidos
            \item Sistemas de medición distribuido y de instrumentación 
        \end{itemize} 
        \item Sistemas de Diálogo Hablado y Multimodal
        \begin{itemize}[nosep]
            \item Speech recognition, understanding and generation
            \item Dialogue management
            \item User simulation
            \item Affective computing
            \item Multimodal interaction
            \item Ambient Intelligence (AmI)
            \item Generation of computer personality models
        \end{itemize}
        \item Tecnologías y Aplicaciones de Realidad Virtual, Interacción y Simulación (TARVIS)
        \begin{itemize}[nosep]
            \item 3D Scanning and dissemination of Cultural Heritage
            \item Numeric simulation in high performance graphic devices.
            \item New technics on Human-Computer Interaction
            \item Solid and Volume Modelling
            \item Medical Applications
        \end{itemize}
    \end{enumerate}
    }
}



\end{document}